% Compile with: uplatex file.tex && dvipdfmx file.dvi
%           or: latexmk -uplatex file.tex
\documentclass[dvipdfmx]{article}
\usepackage{amsmath, amssymb}
\usepackage{tikz}

\title{scenario\_squeezed\_state}
\date{\today}

\begin{document}
\maketitle

A single-mode squeezed state. $a$ is the photon annihilation operator; $b$ is the mode after the squeeze transformation (declared as a generic boson operator, since it is not a named particle species). $r$ is the squeeze parameter.

a as the annihilation operator for photon

\[
  a
\]

b as a boson operator

\[
  b
\]

Define r as a scalar

\[
  r
\]

Register the Bogoliubov transformation $b=S^\dagger a S$ induced by the squeeze operator $S(r)=\exp[\frac{r}{2}(a^2-a^{\dagger 2})]$, both components together (the dagger is not auto-generated, so the creation-operator equation is written out too).

$b(1) = cosh(r) a(1) - sinh(r) a'(1)$

\begin{equation}
  b(1) = \cosh(r) a(1) - \sinh(r) a^{\dagger}(1)
  \label{eq:E1}
\end{equation}

$b'(1) = cosh(r) a'(1) - sinh(r) a(1)$

\begin{equation}
  b^{\dagger}(1) = \cosh(r) a^{\dagger}(1) - \sinh(r) a(1)
  \label{eq:E2}
\end{equation}

Is the transformation physically allowed — does it preserve the canonical commutator $[b,b^\dagger]=1$?

commutator of E1 and E2

\begin{equation}
  \left(\cosh^{2}\left( r \right) - \sinh^{2}\left( r \right)\right)
  \label{eq:E3}
\end{equation}

Express the original mode $a, a^\dagger$ in terms of the squeezed mode $b, b^\dagger$ (the inverse transform).

Invert E1, E2 for a(1), a'(1)

\begin{equation}
  a_{1} = \frac{b^{\dagger}_{1} \cdot \sinh\left( r \right) + b_{1} \cdot \cosh\left( r \right)}{\cosh^{2}\left( r \right) - \sinh^{2}\left( r \right)}
  \label{eq:E4}
\end{equation}

\begin{equation}
  a^{\dagger}_{1} = \frac{b^{\dagger}_{1} \cdot \cosh\left( r \right) + b_{1} \cdot \sinh\left( r \right)}{\cosh^{2}\left( r \right) - \sinh^{2}\left( r \right)}
  \label{eq:E5}
\end{equation}

Simplifying that commutator with the hyperbolic identity gives 1. The transformation is canonical.

simplify E3

\begin{equation}
  1
  \label{eq:E6}
\end{equation}

Register the photon number operator $a^\dagger a$.

a'(1) a(1)

\begin{equation}
  a^{\dagger}(1) a(1)
  \label{eq:E7}
\end{equation}

Substitute the inverse transform, expressing the number operator in the squeezed mode.

In E7, substitute E4 for a(1) and E5 for a'(1)

\begin{equation}
  a^{\dagger}(1) a(1) = \frac{\sinh^{2}\left( r \right)}{\left( \cosh^{2}\left( r \right) - \sinh^{2}\left( r \right) \right)^{2}}  + \frac{\cosh\left( r \right) \sinh\left( r \right)}{\left( \cosh^{2}\left( r \right) - \sinh^{2}\left( r \right) \right)^{2}} {{b}_{1}^{\dagger}}^{2} + \left(\frac{\cosh^{2}\left( r \right)}{\left( \cosh^{2}\left( r \right) - \sinh^{2}\left( r \right) \right)^{2}} + \frac{\sinh^{2}\left( r \right)}{\left( \cosh^{2}\left( r \right) - \sinh^{2}\left( r \right) \right)^{2}}\right) {b}_{1}^{\dagger} {b}_{1} + \frac{\cosh\left( r \right) \sinh\left( r \right)}{\left( \cosh^{2}\left( r \right) - \sinh^{2}\left( r \right) \right)^{2}} {{b}_{1}}^{2}
  \label{eq:E8}
\end{equation}

Take the expectation value in the squeezed vacuum (the $b$ vacuum). The photon number is not zero even in a vacuum — that is what squeezing is.

\begin{equation}
  a^{\dagger}(1) a(1) = \frac{\sinh^{2}\left( r \right)}{\left( \cosh^{2}\left( r \right) - \sinh^{2}\left( r \right) \right)^{2}}
  \label{eq:E9}
\end{equation}

\end{document}